% TOP_PROBLEM: {"id": 4600052, "title": "Flow equivalence of subshifts", "category_id": 11, "category": {"id": 11, "name": "dynamical_systems", "display_name": "Dynamical Systems", "description": "Problems about long-term behavior of deterministic systems, Hamiltonian dynamics, and periodic orbits.", "slug": "dynamical-systems", "order_index": 11, "created_at": "2026-07-31T15:26:25.671Z"}, "status": "partially_solved", "problem_number": "AMR-045-0052", "set_id": 13}
% FIRST_POSTED: 2026-10-01
% ENTRY_KIND: statement_only
% UnsolvedMath ID 4600052; source catalogue code AMR-045-0052
% !TeX program = lualatex
% Definition and sources only; this file is not an LLM solution attempt.
\documentclass[11pt,a4paper]{article}
\usepackage[margin=25mm]{geometry}
\usepackage{fontspec}
\setmainfont{lmroman10-regular.otf}[BoldFont=lmroman10-bold.otf,
 ItalicFont=lmroman10-italic.otf,BoldItalicFont=lmroman10-bolditalic.otf]
\usepackage{amsmath,amssymb,mathtools}
\usepackage{unicode-math}
\setmathfont{latinmodern-math.otf}
\AtBeginDocument{\let\setminus\smallsetminus}
\usepackage{xurl,hyperref}
\hypersetup{colorlinks=true,urlcolor=blue}
\setcounter{secnumdepth}{0}
\setlength{\parindent}{0pt}
\setlength{\parskip}{5pt}
\setlength{\emergencystretch}{3em}
\begin{document}
\section*{Flow equivalence of subshifts}\label{4600052}
{\small UnsolvedMath ID 4600052 · AMR-045-0052 · Dynamical Systems · AMR Open Problem Lists}
\subsection{Definitions and mathematical statement}
Let $(X,\sigma)$ be a nonempty two-sided subshift over a
finite alphabet. Its unit suspension is the compact space
\[
 SX=(X\times\mathbb R)/\!\sim,\qquad
        (x,t+1)\sim(\sigma x,t),
\]
with suspension flow $\phi_X^s[x,t]=[x,t+s]$.
Two subshifts $X,Y$ are flow equivalent if there is a
homeomorphism $H:SX\to SY$ carrying every flow orbit
onto a flow orbit and preserving its time orientation.
Equivalently, along each orbit $H$ admits an increasing
continuous reparametrization of time; preservation of the
actual elapsed time is not required.

Classify subshifts under this relation: give complete
invariants, or a usable necessary and sufficient criterion,
for $X$ and $Y$ to be flow equivalent. The question
includes arbitrary finite-alphabet subshifts, not only those
with finite forbidden-word lists or finite labelled-graph
presentations.

The homeomorphism must be defined on the entire compact
suspension, including closed flow orbits coming from periodic
shift configurations and all their accumulation points.
On a closed orbit it preserves the circle's orientation;
on a nonclosed orbit it preserves the forward direction.
It need not send the chosen cross-section $X\times\{0\}$
to $Y\times\{0\}$. This permits changes in return times
which ordinary topological conjugacy forbids. The problem
asks for a classification of this precise orbit-and-orientation
relation, without imposing a finite input encoding for
all subshifts.
\subsection{Short English statement}
Classify finite-alphabet subshifts when their suspension flows are identified by orientation-preserving orbit homeomorphisms.
\subsection{Sources}
\begin{itemize}
\item[S1] ULAM.AI and the UnsolvedMath Contributors, supplied problems.json, record 4600052 (AMR-045-0052), AMR Open Problem Lists. Catalogue identity and source transcription; CC BY 4.0.
\url{https://huggingface.co/datasets/ulamai/UnsolvedMath}.
\item[S2] Boyle - Open Problems in Symbolic Dynamics (2008), Problem/Question/Conjecture 33.3. Original source indexed by AMR-045-0052.
\url{https://www.math.umd.edu/~mboyle/papers/openfinalsub3nov2007.pdf}.
\end{itemize}
\end{document}
